% Tradução do resumo (capitulos/00-resumo.tex) feita por Claude Code em 18/11/2026.
% Pendente de revisão do autor. Ao alterar o resumo, atualizar este arquivo.
Prompt engineering is one of the ways of using large language models in software development, and studies on the
subject vary widely in techniques, models, tasks and metrics. This work organizes the evidence on prompting
techniques applied to code generation, translation, refactoring and debugging, by means of a rapid review of the
literature, a taxonomy and a decision-support table. The search in arXiv, IEEE Xplore and the ACM Digital Library
returned 2,304 unique studies. Of these, 715 named the task in the title and were screened by two language
models, with checking by the author. At the full-text stage, the author checked the 419 studies read by the models
and included 332. The records of the 332 studies, with 805 comparisons and a quality assessment, were extracted by
AI; the author checked a sample of 67 studies, with corrections in 18 of 571 blocks, and checked against the
original papers the 120 empirical claims made in the text. Most of the evidence compares complete techniques with
baselines, without isolating the mechanism responsible. In the extracted comparisons, many favourable results
involve verifiable external information, such as tests, execution and repository context, and depend on more calls
to the model or on information from the reference solution. Reasoning instructions, self-critique without
execution, fixed examples and changes in wording had mixed results, especially with small models. Fewer than half
of the studies evaluate any risk mitigation. The main limitations are the requirement that the task appear in the
title, the exclusion of studies available only on arXiv and the extraction by AI, checked by the author only in a
sample.
